\section{Diskussion}

In diesem Kapitel werden die wichtigsten Ergebnisse dieser Arbeit kritisch betrachtet und eingeordnet. Ziel der Diskussion ist es, die Bedeutung der Ergebnisse im Hinblick auf die Forschungsfragen (siehe Einführung) herauszuarbeiten, Grenzen zu beleuchten und mögliche Implikationen oder Verbesserungsvorschläge für zukünftige Arbeiten aufzuzeigen. Dabei werden sowohl die Auswahl und Verarbeitung der Merkmale als auch die Leistungsfähigkeit der generierten Modelle analysiert.

\subsection{Merkmalsauswahl}

In diesem Abschnitt werden die Ergebnisse der Merkmalsextraktion und -auswahl diskutiert. Zusammenfassend wurden jeweils elf Merkmale aus drei Prozessierstufen des Datensatzes extrahiert. Die erste Merkmalsgruppe waren die Amplitudenmerkmale, mit der Instantaneous Amplitude, RMS-Amplitude, der Average Energy und der Quadrature. Alle haben besonders die starken Reflexionen der Luft- und Bodenwelle (Bereich \textbf{(A)}) und der oberen Untergrundstrukturen (Bereich \textbf{(B)}) hervorgehoben. Die Average Energy bietet dabei einen höheren Kontrast, da sie in ihrer Berechnung die Amplitudenwerte nochmal quadriert, verliert dadurch aber auch Schärfe, wenn es um schwache Reflektoren geht. Im Gegensatz dazu sind die schwächeren Reflektoren in der Quadrature am besten zu sehen, was besonders an der Struktur \textbf{(C)} sichtbar wird, aber dafür treten auch andere Artefakte, wie die horizontalen Strukturen \textbf{(D1 - D3)}, in den Vordergrund. Diese sind höchst wahrscheinlich keine eigenen Untergrundstrukturen, da sie über 180 m komplett horizontal verlaufen. Außerdem ist zu diesen Laufzeiten ein langwelliges Signal auf allen Radargrammen der Untersuchungsgebietes zu sehen, auf welches die Quadrature, aber auch andere Merkmale sensitiv sind. Aus diesen Erkenntnissen wurde später die dritte Prozessierung durchgeführt, in dem ein Wellenzahlhochpassfilter, dieses Signal entfernen sollte (siehe unten).\\
Zusammenfassend eignen sich alle vier Merkmale sehr gut für eine spätere Segmentierung mittels der SOM, womit die Auswahl eher davon abhängt, auf welche Aspekte das Hauptaugenmerk gelegt werden soll. Möchte man eher die starken Reflektoren hervorheben, sollte man die Average Energy nutzen, für eine glatte Ansicht der Reflektoren eher die RMS Amplitude. In dieser Arbeit haben sich die Instantaneous Amplitude in Prozessierstufe 2 und die Quadrature in Prozessierstufe 3 am besten nachbearbeiten lassen und zeigten im Ergebnisplot die erwarteten Ergebnisse. Wichtig ist, dass pro SOM-Training nur eines der Merkmale verwendet wird, da diese eine hohe Korrelation miteinander aufweisen.\\
Die Instantaneous Phase ist ein solides Merkmal, welches ohne Nachbearbeitung auskommt und besonders gut die Untergrundstrukturen sichtbar macht, indem sie die Laufzeitunterschiede der Reflexionen hervorhebt. Da eine Self-Organizing Map allerdings versucht gleiche Eingabewerte zu clustern, eignet sie sich für diese Art der Segmentierung nicht, da sie als periodische Größe ständig ihre Werte mit der Laufzeit ändert. Dies ist auch an den Ergebnissen des Modells 3 (siehe Abb. \ref{fig:model_3_bmu}) zu sehen, indem die Untergrundstrukturen in Bereich \textbf{(A)} aufgrund der Phase zu verschiedenen Neuronen zugeordnet wurden. Interessant sind für diese Arbeit die Änderungsraten der Phase, welche die Instantaneous Frequency als zeitliche Ableitung der Instantaneous Phase darstellt.\\
Die Instantaneous Frequency zeigt die Momentanfrequenz des Signals an, welches bei einer Antennenmittenfreqenz von 200 MHz, auch die Untergrundstrukturen \textbf{(A)}, in einem Bereich um diese Frequenz darstellt. Bereiche mit wenig Reflektivität besitzen geringere Frequenzen, weshalb sich dieses Merkmal sehr gut für die Segmentierung eignet. Allerdings zeigt die Frequenz auch die Schwierigkeiten dieses Datensatzes auf, da periodische Störsignale in Zeit- \textbf{(A - C)} sowie Wegrichtung \textbf{(D)} sehr gut zu erkennen sind, welche sich auch auf allen anderen Radargrammen wiederfinden. In der Prozessierstufe 3 mit der stärkeren Nachbearbeitung, sind diese Artefakte zwar nicht mehr zu sehen, aber dafür entfällt auch die klare Trennung zwischen strukturierten Bereich und Bereichen ohne Struktur (siehe Abb. \ref{fig:inst_freq_p3}), weshalb in dieser Arbeit nur die Prozessierstufe 1 verwendet wurde.\\
Das letzte Momentanattribut der Instantaneous Quality Factor oder auch kurz Instantaneous Q, gibt die Güte des Signals an. Dieser zeigt ebenfalls hohe Werte in Bereichen mit hoher Struktur (\textbf{(A)}, \textbf{(C)}), besitzt aber keine klare Trennung zwischen Bereichen hoher und niedriger Struktur, wie zum Beispiel die Amplitudenattribute, welche die gleichen Bereiche hervorheben. Desweiteren basiert seine Berechung direkt auf der Instantaneous Amplitude und der Instantaneous Frequency, weshalb er mit ihnen teilweise korreliert. Da der Q-Faktor allerdings niedrige Werte in Bereich \textbf{(B)} besitzt, welcher von einem Störsignal geprägt ist, eignet er sich teilweise für eine Segmentierung, wahrscheinlich mit einer besseren Nachbearbeitung. In dieser Arbeit wurde er allerdings nicht weiter untersucht, da die anderen Merkmale bereits eine gute Segmentierung ermöglichten.\\
Die Statistischen Momente Skewness und Kurtosis sind zwei Merkmale, die sich sehr stark ähneln und wahrscheinlich austauschbar genutzt werden können. In dieser Arbeit wurde ausschließlich die Kurtosis verwendet, da sie visuell einen höheren Konstrast aufweist. Sie zeigt mit hohen Werten sehr eindeutig Bereiche mit Rauschen und hat dort negative Werte, wo das Signal klar ist, vor allem, wenn es noch stärkere Reflexionen zu längeren Laufzeiten gibt (\textbf{(A)}, \textbf{(C)}). Da das Störsignal bei ca. 20 m \textbf{(B)} ein geringes Rauschen besitzt, können die statistischen Momente dieses gut detektieren und von der Umgebung abheben. Des Weiteren ist zwischen Prozessierstufe 2 (Abb. \ref{fig:kurtosis_p2}) und 3 kein Unterschied zu sehen, was nach \textcite{barnes_too_2006} für eine Robustheit gegen Änderungen spricht. Allgemein eignen sich diese Merkmale sehr gut für eine Segmentierung in dieser Arbeit, weshalb sie in jedem Lauf verwendet wurden. Wichtig ist nur, dass nur eines der beiden Merkmale verwendet wird, damit die Unabhängigkeit von allen anderen Merkmalen in einem Trainingslauf erhalten bleibt.\\
Die Semblance eignet sich als Änhlichkeitsattribut nur teilweise für eine Segmentierung: Sie zeigt in Prozessierstufe 2 zwar Bereiche hoher Ähnlichkeit \textbf{(A)} dort, wo auch andere Merkmale wenige Strukturen darstellen und niedrige Ähnlichkeit \textbf{(B)} in Bereichen mit viel Struktur, ist aber sehr sensitiv auf die langwelligen Störsignale in Zeitrichtung \textbf{(C)}. In Prozessierstufe 3 lieferte sie nicht annähernd gleiche Ergebnisse (siehe Abb. \ref{fig:semblance_p3}), weshalb sie in dieser Arbeit im Modell 4 nicht weiter verwendet wurde.\\
Der Dip ist ebenfalls ein Merkmal, welches in Prozessierstufe 2 zwar Bereiche mit klarer Neigung dort zeigt, wo es klare Untergrundstrukturen gibt \textbf{(A)}, aber in Bereichen mit wenig Struktur \textbf{(B)} eine zu starke Varianz aufweist. In der Form, wie sie in dieser Arbeit berechnet wurde, eignet sie sich für eine Segmentierung nicht. Da es aber noch einige andere Berechnungsarten gibt, ist nicht auszuschließen, dass ggf. bessere Ergebnisse erzielt werden können und so auch Bereiche verschiedenen Dips im Untergrund in verschiedene Segmente unterteilt werden können.\\
Zusammenfassend lässt sich sagen, dass sich die Merkmale Instantaneous Amplitude, Quadrature, Instantaneous Frequency und Kurtosis am besten geeignet haben, um ein Radargramm zu segmentieren.

\subsection{Modellvergleich}

Zu Beginn dieser Arbeit stellte sich die Frage, ob und wie gut eine Self-Organizing Map aus den berechneten Merkmalen sinnvolle Cluster erzeugen kann. Dafür wird zunächst geschaut, wie viele Epochen notwendig sind, um ein gutes Modell zu trainieren. Der Epochenvergleich nach $N = 10$, $30$ und $50$ Epochen für Modell 1 hat gezeigt, dass bereits nach 10 Epochen ein gutes Ergebnis erzielt wurde (siehe Abb. \ref{fig:model_1_planes_10}). Die Ergebnisbilder für die einzelnen Durchläufe sind nur schwer voneinander zu unterscheiden. Einzig die U-Matrix zeigt, dass die Neuronen nach 10 und 30 Epochen weniger nah beieinander liegen, also weniger dichte Cluster bilden. Des Weiteren ist die Wahl der Hyperparameter entscheidend für die Epochenanzahl, da eine kleinere Startlernrate auch eine höhere Epochenanzahl benötigt um sich ähnlich zu sortieren. In diesem Fall haben sich die Faustformeln und Standardgrößen für die Wahl der Kartengröße, der Startlernrate, der Lernratenabnahmefunktion, des Nachbarschaftsradius und der Abnahmefunktion des Nachbarschaftsradius bewährt.\\
Im Vergleich von Modell 1 mit Modell 2 zeigte sich, dass für diese Eingabedaten die Hyperparameter $\alpha_0 = 0.3$ und $\sigma_0 = 7$ zwar bei der Abbildung der Daten änhlich gut funktionieren, wie $\alpha_0 = 0.5$ und $\sigma_0 = 5$, da der $QE$ ungefähr gleich war, aber die Topologieerhaltung mit einem größeren Nachbarschaftsradius besser funktioniert (siehe Tabelle \ref{tab:model_error_overview}). Daher wurde diese Kombination auch für alle weiteren Läufe beibehalten.\\
Interessanterweise zeigt Modell 3 durch die Merkmalauswahl deutlich schlechtere Ergebnisse im Bezug auf den Quantisierungsfehler, aber besitzt den besten Topologieerhalt mit einem $TE$ von rund $0,047$. Modell 4 hingegen hat den besten Quantisierungsfehler, aber dafür den schlechtesten Topologieerhalt. In Tabelle \ref{tab:model_error_overview} sind dafür zusätzlich die Anzahl der verwendeten Merkmale angegeben. Vergleicht man die Fehlergrößen mit der Anzahl der Merkmale $K$ bei 50 Epochen, so zeigt sich eine Korrelation von $QE$ und $K$ und eine Antikorrelation von $TE$ und $K$. Da es hierzu bisher keine Forschungsergebnisse gibt, lässt sich nicht endgültig sagen, ob dies allgemein der Fall ist. Dies gilt besonders, da unterschiedliche Merkmale verwendet wurden und der Vergleich in dieser Arbeit auf nur vier Modelle limitiert ist. Allerdings scheint es logisch, dass bei gleicher Kartengröße und wenigen Merkmalen die Clustergröße im Sinne der Neuronenzahl pro Cluster zunimmt und entsprechend bei mehr unabhängigen Merkmalen auch mehrere Cluster mit weniger Neuronen pro Cluster entstehen. Da die Berechnung des $TE$ davon abhängt, ob die SBMU ein direkter Nachbar der BMU ist oder nicht, erscheint es ebenfalls logisch, dass die Wahrscheinlichkeit, dass dies nicht der Fall ist, bei dichteren Clustern mit mehr Neuronen höher ist. Gleiches gilt umgekehrt für den Quantisierungsfehler. Eine höhere Neuronendichte im Cluster würde logischerweise zu geringeren Abständen zwischen den Eingabepunkten zu ihrer BMU führen. Dies wäre ein Sachverhalt, welcher in Folgestudien genauer untersucht werden könnte.\\
Die Trefferkarte von Modell 4 (Abb. \ref{fig:model_4_hits}) zeigt, dass die Neuronen mit den größten Trefferzahlen alle am Rand liegen. Dies ist bei dem Großteil der anderen Modelle auch der Fall, allerdings mit einer weniger starken Ausprägung. Dieser sogenannte \textit{Edge Effect} wird schon von \textcite{schmidt_effects_2011} beschrieben und ist ein bekanntes Phänomen, besonders in der Anwendung von traditionellen zweidimensionalen quadratischen Neuronengittern. Werden - wie in dieser Arbeit - 2D Colormaps verwendet, so sollte dieser Effekt unbedingt bedacht werden (siehe \textcite{steiger_explorative_2015}), da  Colormaps , welche höhere Konstraste am Rand oder höhere Konstraste in der Mitte besitzen, im Ergebnisplot unterschiedlich gute Ergebnisse erzeugen können. Des Weiteren hat die Clustergröße hier einen direkten Einfluss auf die Abgrenzung der Farben im Ergebnisplot. Besitzen Cluster eine hohe Dichte, also eine hohe Neuronenzahl, so ist die Farbvarianz im Ergebnisplot ebenfalls groß. Dies ist z.B. im Ergebnisplot (Abb. \ref{fig:model_4_bmu}) zu sehen. Dort erfahren Bereiche im unteren Teil des Bildes einen Farbverlauf von grün über gelb zu rot-orange, obwohl die Neuronen am linken Rand des Gitters (siehe Abb. \ref{fig:model_4_u_matrix}) nah beieinander liegen. Dass die Darstellung in diesem Bereich von einem verwaschenen grün zu einem rot mit scharfen Kanten wechselt, liegt aber auch an der Nachbearbeitung der Merkmale. Während die Instantaneous Frequency in einem 21 x 21 Fenster geglättet wurde (siehe Tabelle \ref{tab:feature_params_1}), wurde die Quadrature in einem fast vier Mal größerem Fenster (81 x 21, siehe Tabelle \ref{tab:feature_params_3}) deutlich stärker geglättet. Dieser Unterschied ist nach dem Training auch im Ergebnisplot zu sehen. Für ein gleichmäßigeres Ergebnisbild sollte in zukünftigen Arbeiten darauf geachtet werden, dass die Glättung der Merkmale einheitlich ist.\\
Allgemein lässt sich aber sagen, dass der Ergebnisplot von Modell 4 (Abb. \ref{fig:model_4_bmu}) die Untergrundstrukturen in blau-lila gut von den eher homogeneren Bereichen abgrenzt, wobei blaue Bereiche stärkere Reflexionen und lila Bereiche schwächere Reflexionen bedeuten. Die rot-orangen Cluster zeigen Bereiche geringer Signalstärke und verhältnismäßig stärkeren Rauschens, während die grün-gelben Bereiche zu kürzeren Laufzeiten homogene Bereiche mit weniger starkem Rauschen darstellen. Das Modell zeigt ein Störsignal bei ungefähr 20 m Laufzeit, welches ebenfalls im Radargramm (siehe Abb. \ref{fig:radargramm_1}) zu sehen ist. Zusammenfassend erzeugt dieses SOM-Modell also sinnvolle Cluster, welche die Untergrundstrukturen gut voneinander abgrenzen und auch Störsignale detektieren können.\\